\section{Minimum Multiplications to Reach End}
\label{sec:graph-min-multiplications}

\problemstatement{Given a start number, an end number, and an array of
multipliers, in each step you may replace the current number $x$ with
$(x \times \textit{arr}[i]) \bmod 100000$ for any $i$. Find the minimum
number of steps to turn \textit{start} into \textit{end}, or $-1$ if
impossible.}

\begin{patternbox}
The disguise here is strong: there is no explicit graph. But each
reachable number $0\ldots99999$ is a \textbf{node}, and each multiplier
defines an \textbf{edge} of cost one step. ``Minimum steps, all costs
equal'' is again a \textbf{unit-weight shortest path}, so \textbf{BFS}
over the $100000$ possible numbers solves it.
\end{patternbox}

\subsection*{Modeling numbers as nodes}

The modulo keeps every reachable value in $[0,99999]$, a finite state
space. From a value $x$, applying multiplier $\textit{arr}[i]$ produces
a neighbor $(x\cdot\textit{arr}[i])\bmod 100000$. Every such transition
is one step, so all edges have equal weight and BFS from \textit{start}
finds the fewest steps to \textit{end}. A visited array over the
$100000$ states prevents revisiting and guarantees termination.

\begin{ideabox}[A Bounded State Space Makes BFS Finite]
The $\bmod\,100000$ is what makes this tractable: without it the numbers
could grow without bound. With it, there are only $100000$ possible
states, so BFS explores each at most once and either reaches
\textit{end} or exhausts the reachable set.
\end{ideabox}

\subsection*{Algorithm}

\begin{enumerate}
  \item If \textit{start} equals \textit{end}, return $0$. Push
        \textit{start} with $\textit{steps}=0$; mark it visited.
  \item While the queue is non-empty, pop $(\textit{value},
        \textit{steps})$.
  \item For each multiplier, compute
        $\textit{next}=(\textit{value}\cdot\textit{arr}[i])\bmod
        100000$; if $\textit{next}=\textit{end}$, return
        $\textit{steps}+1$; if unvisited, mark and push it with
        $\textit{steps}+1$.
  \item If the queue empties, return $-1$.
\end{enumerate}

\subsection*{C++ implementation}

\begin{lstlisting}
#include <bits/stdc++.h>
using namespace std;

int minimumMultiplications(vector<int>& arr, int start, int end) {
    if (start == end) return 0;
    const int MOD = 100000;
    vector<bool> vis(MOD, false);
    queue<pair<int,int>> q;                  // {value, steps}
    q.push({start, 0});
    vis[start] = true;

    while (!q.empty()) {
        auto [value, steps] = q.front(); q.pop();
        for (int m : arr) {
            int next = (int)((long long)value * m % MOD);
            if (next == end) return steps + 1;
            if (!vis[next]) {
                vis[next] = true;
                q.push({next, steps + 1});
            }
        }
    }
    return -1;
}

int main() {
    vector<int> arr = {2, 5, 7};
    cout << minimumMultiplications(arr, 3, 30) << endl;   // 2  (3*2=6, 6*5=30)
    return 0;
}
\end{lstlisting}

\subsection*{Dry run}

\dryrun{\textit{arr}$=\{2,5,7\}$, start $3$, end $30$.}

\begin{itemize}
  \item Step $0$: value $3$. Neighbors: $6$ (via $2$), $15$ (via $5$),
        $21$ (via $7$) -- none is $30$; enqueue all with steps $1$.
  \item Step $1$, value $6$: $6\times2=12$, $6\times5=30$ --
        \textbf{match!} Return $1+1=2$.
  \item Shortest: $3\xrightarrow{\times2}6\xrightarrow{\times5}30$, two
        steps.
\end{itemize}

\subsection*{Complexity}

\begin{complexitybox}
\textbf{Time Complexity.} $O(100000 \times \lvert\textit{arr}\rvert)$ --
each of the $100000$ states is expanded once, trying every multiplier.
\textbf{Space Complexity.} $O(100000)$ for the visited array and queue.
\end{complexitybox}

\begin{takeawaybox}
When a problem asks for the minimum number of equal-cost transformations
between two values, model each value as a node and each transformation
as a unit edge, then run BFS -- even when no graph is stated. A bounded
state space (here, the modulo) is the signal that this modeling is
finite and BFS applies.
\end{takeawaybox}
